\section*{अध्याय \ref{ch:b3:electronic-spectroscopy} --- इलेक्ट्रॉनिक स्पेक्ट्रमिकी और प्रकाशभौतिकी}

\begin{solution}{exo:b3:electronic-spectroscopy:1}
(a) चक्रण नियम द्वारा वर्जित ($\Delta S = 1$)। (b) \omterm{thm:b3:electronic-spectroscopy:laporte}{लापोर्ट नियम} द्वारा वर्जित
($g \to g$); कंपन-इलेक्ट्रॉनिक युग्मन के कारण दुर्बल रूप में दिखता है। (c) अनुमत: $B_{1u}$,
$z$ का निरूपण है, $D_{2h}$ में। (d) सममिति द्वारा वर्जित: $A_2$, $x$,
$y$, $z$ में से कोई नहीं है, $C_{2v}$ में; बैंड दुर्बल है।
\end{solution}

\begin{solution}{exo:b3:electronic-spectroscopy:2}
$10^7/349 - 10^7/450 = 28653 - 22222 = \qty{6431}{cm^{-1}}$, \qty{0.80}{eV}।
\end{solution}

\begin{solution}{exo:b3:electronic-spectroscopy:3}
$A = 5700 \times 1 \times 1.0\times10^{-5} = 0.057$; अवशोषित अंश $1 - 10^{-0.057} =
0.12$।
\end{solution}

\begin{solution}{exo:b3:electronic-spectroscopy:4}
$k_r = \Phi_F/\tau = \qty{1.5e8}{s^{-1}}$, $\tau_r = \qty{6.7}{ns}$; $k_{nr} = (1 - \Phi_F)/\tau
= \qty{1.0e8}{s^{-1}}$।
\end{solution}

\begin{solution}{exo:b3:electronic-spectroscopy:5}
$\int\varepsilon\,\dd\tilde\nu = 1.0645 \times 5700 \times 5000 = \num{3.03e7}$, अतः $f \approx
4.32\times10^{-9} \times 3.03\times10^7 = 0.13$: मध्यम सामर्थ्य का एक अनुमत
संक्रमण।
\end{solution}

\begin{solution}{exo:b3:electronic-spectroscopy:6}
प्वासों गुणक $\eu^{-S}S^v/v!$, $v = \lfloor S\rfloor$ पर अधिकतम है। $S = 0.3$: स्वयं
0--0 रेखा। $S = 1.5$: $v' = 1$, 0--0 रेखा का 1.5 गुना। $S = 5$: $v' = 4$
और 5 समान रूप से, 0--0 रेखा का 26 गुना।
\end{solution}

\begin{solution}{exo:b3:electronic-spectroscopy:7}
$\sum k = \qty{4.2e8}{s^{-1}}$: $\tau = \qty{2.4}{ns}$, $\Phi_F = 0.24$, $\Phi_{\mathrm{ISC}} = 0.71$
(और $\Phi_{\mathrm{IC}} = 0.05$)।
\end{solution}

\begin{solution}{exo:b3:electronic-spectroscopy:8}
$\Phi = 0.546 \times 0.46 \times (0.050/0.040) \times (1.4266/1.333)^2 = 0.36$।
\end{solution}

\begin{solution}{exo:b3:electronic-spectroscopy:9}
$\tau_0/\tau = 1.429$ और 1.860: $[Q]$ में रैखिक, अतः स्वयं आयुकाल
घटता है: \omterm{def:b3:electronic-spectroscopy:quencher}{गतिक शमन}। $K_{SV} = 43$ L/mol (दोनों बिंदु), $k_q = 43/8.0
\times 10^{-9} = \qty{5.4e9}{L.mol^{-1}.s^{-1}}$।
\end{solution}

\begin{solution}{exo:b3:electronic-spectroscopy:10}
\omterm{def:b3:electronic-spectroscopy:quencher}{स्थैतिक शमन}: जो अणु अब भी उत्सर्जन करते हैं वे पहले जितना ही जीवित रहते हैं; उनमें से
आधे मूल अवस्था में एक अप्रतिदीप्त संकुल में बँधे हैं। संगुणन स्थिरांक
$K_a$ के साथ मुक्त अंश $1/(1 + K_a[Q])$ है और $I_0/I = 1 +
K_a[Q]$, जो स्टर्न--वोल्मर जैसा ही रूप है परंतु $\tau$ अपरिवर्तित रहता है।
\end{solution}

\begin{solution}{exo:b3:electronic-spectroscopy:11}
$\tau_r = \tau/\Phi_F = \qty{14}{ns}$; $\Phi_T = 1 - 0.36 = 0.64$। मापा गया आयुकाल
प्रतिस्पर्धी \omterm{def:b3:electronic-spectroscopy:radiationless}{अंतरतंत्र संक्रमण} से घट जाता है; $\tau_r$ वह है जो केवल उत्सर्जन
देता।
\end{solution}

\begin{solution}{exo:b3:electronic-spectroscopy:12}
$\frac12\langle\alpha\beta - \beta\alpha|\alpha\beta + \beta\alpha\rangle = \frac12(1 + 0 - 0 - 1) = 0$, हर इलेक्ट्रॉन के लिए
$\langle\alpha|\alpha\rangle = \langle\beta|\beta\rangle = 1$, $\langle\alpha|\beta\rangle = 0$ का उपयोग करके।
द्विध्रुव \omterm{def:b3:quantum-model-systems:operator}{संकारक} चक्रण को नहीं छूता, अतः संक्रमण आघूर्ण में यह
शून्य गुणक आता है: $T_1 \leftarrow S_0$ की कोई तीव्रता नहीं होती, जब तक \omterm{def:b3:many-electron-atoms:spin-orbit}{चक्रण--कक्षक युग्मन}
अवस्थाओं को मिश्रित न करे।
\end{solution}

\begin{solution}{pb:b3:electronic-spectroscopy:1}
\textbf{1.} $A = 5700 \times 1 \times 1.0\times10^{-4} = 0.57$।
\textbf{2.} $1 - 10^{-0.57} = 0.73$।
\textbf{3.} कुछ हज़ार का $\varepsilon$: मध्यम सामर्थ्य का एक अनुमत $\pi\to\pi^*$
संक्रमण।
\textbf{4.} $1239.8/349 = \qty{3.55}{eV}$।
\textbf{5.} वह केवल पराबैंगनी में अवशोषण करती है; दृश्य प्रकाश पार हो जाता है।
\textbf{6.} ताकि अवशोषित प्रकाश सांद्रता के समानुपाती रहे और उत्सर्जित प्रकाश के
पुनः-अवशोषण (आंतरिक-फ़िल्टर प्रभाव) से बचा जा सके।
\textbf{7.} $S_0 \to S_1$ (अवशोषण, फिर \omterm{def:b3:electronic-spectroscopy:radiationless}{कंपनिक विश्रांति}); $S_1$ से:
\omterm{def:b3:electronic-spectroscopy:luminescence}{प्रतिदीप्ति}, \omterm{def:b3:electronic-spectroscopy:radiationless}{आंतरिक रूपांतरण}, $T_1$ में \omterm{def:b3:electronic-spectroscopy:radiationless}{अंतरतंत्र संक्रमण}।
\textbf{8.} $k_r = 0.546/19 \times 10^{-9} = \qty{2.87e7}{s^{-1}}$; $\tau_r = \qty{35}{ns}$।
\textbf{9.} $(1 - 0.546)/\tau_0 = \qty{2.39e7}{s^{-1}}$।
\textbf{10.} $28653 - 22222 = \qty{6431}{cm^{-1}}$।
\textbf{11.} उत्सर्जन विश्रांत $S_1$ से आरंभ होता है और $S_0$ के कंपनिक रूप से उत्तेजित
स्तरों पर समाप्त होता है, अणु और विलायक के विश्रांत होने के बाद: फ़ोटॉन में कम
ऊर्जा होती है, यहाँ इतनी कि वह दृश्य क्षेत्र में आ जाए।
\textbf{12.} $\Phi_F \times (349/450) = 0.42$: अवशोषित ऊर्जा का 42\,\%।
\textbf{13.} बिंदु हर \qty{0.010}{mol/L} पर 0.594 बढ़ते हैं: 1 से होकर जाती एक सीधी रेखा।
\textbf{14.} पाँचों बिंदुओं से होकर न्यूनतम वर्ग: ढलान $K_{SV} = \qty{59.4}{L/mol}$,
अंतःखंड 1.000।
\textbf{15.} $k_q = K_{SV}/\tau_0 = \qty{3.1e9}{L.mol^{-1}.s^{-1}}$।
\textbf{16.} $\tau = \tau_0/(1 + 59.4 \times 0.040) = \qty{5.6}{ns}$।
\textbf{17.} आयुकाल मापिए: \omterm{def:b3:electronic-spectroscopy:quencher}{गतिक शमन} के लिए $\tau_0/\tau$ उसी रेखा पर पड़ता है
जिस पर $I_0/I$।
\textbf{18.} टॉनिक जल में शमन करने वाला कोई क्लोराइड नहीं है, और वह इतना अम्लीय है कि
क्विनीन को प्रोटॉनित, अर्थात् उसके प्रतिदीप्त रूप में, बनाए रखे।
\textbf{19.} $k_d = 8 \times 8.314 \times 298/(3 \times 0.890\times10^{-3}) =
\qty{7.4e6}{m^3.mol^{-1}.s^{-1}} = \qty{7.4e9}{L.mol^{-1}.s^{-1}}$।
\textbf{20.} $k_q$, $k_d$ के आधे से थोड़ा कम है।
\textbf{21.} लगभग 42\,\%।
\textbf{22.} मंद: $k_d \propto 1/\eta$, और $k_q$ उसके साथ केवल घट ही सकता है।
\textbf{23.} \textbf{$k_q/k_d = 0.42$}: क्लोराइड क्विनीन का शमन विसरण-नियंत्रित
भेंटों की लगभग आधी दर से करता है।
\end{solution}
